\scanpage{122}

2) $F$ \deleted{\unsure{$=\mathbb Q_{\ell X}$}} strictement spécial\inserted{ de poids $j$}, $X$ régulier.

On peut supposer $F=R^dp_*(\mathbb Q_{\ell X'})$, \fbox{$p$ de dim relative $d$}

\diagram{\begin{tikzcd}[ampersand replacement=\&,column sep=large,row sep=large,execute at end picture={\draw[rounded corners=2pt] ([xshift=3mm,yshift=5mm]xp.east) rectangle ([xshift=4mm,yshift=-7mm]yp.east);}]
|[alias=xp]|X'\arrow[r,hook,no head,"f'"]\arrow[d,"p"'] \& |[alias=yp]|Y'\arrow[d,"q"]\\
X\arrow[r,hook,"f"'] \& \deleted{\unsure{X}}Y
\end{tikzcd}}

On a alors

\[ \unsure{\mathbb R|R}^*(fp)_*(\mathbb Q_{\ell X'})\Longleftarrow E_2^{pq}=\mathbb R^p f_*(R^q\unsure{g|p|\varphi}_*(\mathbb Q_\ell)) \]

suite spectrale qui \uline{dégénère} par Deligne, donc $R^if_*(F)=R^if_*(R^j g_*(\mathbb Q_{\ell\unsure{X|X'}}))$ est un \uline{morceau} de $R^{i+j}(fp)_*(\mathbb Q_{\ell X'})$

\[ \text{Or}\qquad R^*(fp)_*(\mathbb Q_{\ell X'})\Longleftarrow \deleted{\unsure{E}}E_2^{pq}=R^p q_*(R^q f'_*(\mathbb Q_{\ell X'})) \]

Or $Y'$ est de dim \deleted{\unsure{$n+j$}} $n'\leq\text{\tikz[baseline=(a.base)]\node[draw,ellipse,inner sep=1pt](a){$n+j$};}$

\[
\begin{aligned}
\text{donc}\qquad R^q f'_*(\mathbb Q_{\ell X'})&\quad\text{est effectif de poids }\leq 2\operatorname{Inf}(\deleted{\unsure{2}}q,\deleted{\unsure{2}}n+\deleted{\unsure{2}}j)\\
&\text{de niveau }\leq\operatorname{Inf}(q,n+j-1)
\end{aligned}
\]

$\unsure{q|g}$ est de dim relative $\leq$

\deleted{\unsure{On prend} $p+q=i+j$}

\deleted{\unsure{avec} $p+q\unsure{=}$} $E_2^{pq}$ de poids $\leq$

% archive sheet number: 121
